\begin{tikzpicture}[font=\small,>=stealth]
\draw[rounded corners=3pt,gray!45] (0,0) rectangle (14.2,13.2);
\node at (7.1,12.55) {放回双曲线：同支、反向，两段的倒数相加};\begin{scope}[shift={(4.5,7.5)},scale=0.37]
\begin{scope}\clip (-10,-9) rectangle (10,9);
\draw[teal!75!black,thick,samples=110,domain=-1.9:1.9,variable=\t] plot ({3*cosh(\t)},{4*sinh(\t)});
\draw[teal!75!black,thick,samples=110,domain=-1.9:1.9,variable=\t] plot ({-3*cosh(\t)},{4*sinh(\t)});
\end{scope}
\draw[->,gray] (-10.2,0)--(10.3,0);
\draw[orange!85!black,very thick] (5.0,0)--(4.2000000000,-3.9191835885);\draw[blue!75!black,very thick] (5.0,0)--(6.6000000000,7.8383671769);
\fill (6.6000000000,7.8383671769) circle(6pt) node[above right=7pt] {$A$};\fill (4.2000000000,-3.9191835885) circle(6pt) node[below=8pt] {$B$};
\fill (5.0,0) circle(6pt) node[below right=9pt] {$F$};
\end{scope}

\node[anchor=west,align=left] at (9.5,8.3)
 {图取 $a=3,b=4,c=5$\\[13pt]$AF=8,\quad BF=4$\\[13pt]两点都在右支\\[13pt]$F$ 在 $AB$ 内\\[13pt]方向相反};
\node at (7.1,3.65) {$\displaystyle \frac18+\frac14=\frac38=\frac{2a}{b^2}$};
\fill[blue!30] (3,2.25) rectangle (4.4,2.75);
\fill[orange!35] (4.4,2.25) rectangle (7.2,2.75);
\draw[gray] (3,2.25) rectangle (7.2,2.75);
\node[blue!75!black] at (3.7,1.65) {$1/AF=1/8$};
\node[orange!85!black] at (5.8,1.0) {$1/BF=1/4$};
\node[anchor=west] at (8.1,2.5) {两块合起来，才是固定的 $3/8$};
\node at (7.1,.45) {条形图表示倒数；上面的蓝、橙线段表示原长度，两者不能混读。};
\end{tikzpicture}
